\section{An exact choice of the charges and nuclear positions}
\label{sec:parameters}

We now specify one member of the family without leaving either
large parameter free.  The resulting charge is an exact integer
defined through spectral and variational quantities.
The definition is nonnumerical: no evaluated decimal value or
effective numerical upper bound for this integer is asserted.

Throughout this section fix
\begin{equation}
 D=10^{200},\qquad \eta_0=10^{-60},\qquad
 L=10,\qquad \mu=\frac1{10}.
 \label{eq:fixed-parameter-numbers}
\end{equation}
The orbitals \(b,e\), the numbers \(\epsilon_b,\epsilon_e,\gamma,k_0\),
and the reference angle
\[
 \theta_*=\arctan\frac{e}{\sqrt2\,b},\qquad
 \delta_*=\frac\pi2-\theta_*\quad(z\ge0)
\]
always refer to this fixed two-centre problem at zero coupling.
For auxiliary couplings \(t\in[0,1]\), use the operators
\(H_{m,t}\) of \Cref{eq:scaled-coulomb-hamiltonian}.
Subscripts \(t\) on all other quantities refer to these auxiliary
operators.

\begin{definition}[Regular auxiliary data]
\label{def:regular-data}
An auxiliary coupling \(t\in[0,1]\) is called regular if all the
following conditions hold.
\begin{enumerate}[label=\textup{(\roman*)}]
\item The two-electron ground level is isolated and simple, and its
ground state is an axisymmetric, reflection-even spin singlet.
The three-electron ground level in \(S_z=1/2\) is isolated and
simple, is an absolute ground level, and has an axisymmetric,
reflection-odd ground state.
\item The spin-summed density \(n_t\) is everywhere positive, and
\begin{equation}
 2\epsilon_b-E_3(t)>0.49,\qquad 0.99<k_t<\gamma,
 \qquad k_t=\sqrt{2(E_2(t)-E_3(t))}.
 \label{eq:regular-energy-tests}
\end{equation}
\item With the ionic ground projection \(P_t\) in the
distinguished-electron fiber space, and with
\[
 F_{{\rm true},t}(x)=\sqrt3\,\Psi_t(x,\cdot),\qquad
 A_t(r)=e^{-k_t r}r^{(2-2t)/k_t-1},
\]
the direct-removal amplitude satisfies
\begin{equation}
 \inf_{\omega_z>0}\ \liminf_{r\to\infty}
 \frac{\|P_tF_{{\rm true},t}(r\omega)\|}
      {\omega_z A_t(r)}>0 .
 \label{eq:regular-amplitude-test}
\end{equation}
\end{enumerate}
If any quantity required for these tests is undefined, the coupling
is declared nonregular.
\end{definition}

The simple levels make the densities and projected norms
independent of the phases chosen for their state vectors.
The Hamiltonians commute with complex conjugation, so their simple
eigenvectors may be chosen real and normalized.
For every regular \(t\), \Cref{lem:density-regularity} gives the
Sobolev regularity needed below.  For \(t>0\), smoothness away from
the nuclei follows from \Cref{prop:fixed-parameter-channels}; at
\(t=0\), it follows from \(n_0=2b^2+e^2\) and the scalar orbital
equations for \(b,e\).
In particular, the exterior density has the derivatives occurring
in the next definition.

\subsection{The radius and its strict comparison margin}

For an integer \(j>10D\) and regular \(t\), define
\(m_t(p,z)=e^{2k_t p}n_t(p,z)\), and let \(H_t(j)\) be the maximum of
the following four numbers:
\begin{align}
 &\frac1{\eta_0}
   \sup_{|x|\ge j}
    \left|\partial_r\log n_t(x)+2k_t\right|,
 \notag\\
 &\frac1{\eta_0}
   \sup_{\substack{|x|\ge j\\0\le z\le p/2}}
    \left|\partial_p\log n_t(x)+2k_t p/r\right|,
 \notag\\
 &\frac1{\eta_0}
   \sup_{\substack{p\ge j\\0\le z\le L}}
    \frac{|\partial_pm_t(p,z)|+|\partial_p^2m_t(p,z)|}
         {m_t(p,z)},
 \notag\\
 &\frac1{100}
   \sup_{\substack{p\ge j\\0\le z_1\le z_2\le L}}
    \frac{n_t(p,z_1)}{n_t(p,z_2)}.
 \label{eq:parameter-H}
\end{align}
Here \(\partial_p\) holds \(z\) fixed.
For nonregular \(t\), set \(H_t(j)=+\infty\).

Let \(M\) be the smallest integer in the following set:
\begin{equation}
 \left\{
 \begin{aligned}
 j\in\N:\;&j>10D,\quad
 \max\left\{H_0(j),\,\limsup_{t\downarrow0}H_t(j)\right\}<1,\\
 &\sup_{0\le z\le L}
 \frac{\partial_z\bigl(n_0(j,z)\partial_z\delta_*(j,z)\bigr)}
      {n_0(j,z)\delta_*(j,z)}
 <\mu^2
 \end{aligned}
 \right\}.
 \label{eq:parameter-radius}
\end{equation}
The coordinates \((j,z)\) are cylindrical coordinates.  For regular
\(t\), let \(U_t(p,z)\) be the strip solution
\begin{equation}
 \frac1{n_t(p,z)}
 \partial_z\bigl(n_t(p,z)\partial_zU_t(p,z)\bigr)
 =\mu^2U_t(p,z),\qquad
 U_t(p,0)=1,\quad U_t(p,L)=0 .
 \label{eq:parameter-strip-solution}
\end{equation}
It is the normalized positive solution constructed in
\Cref{lem:strip-solution}.  Once \(M\) is defined, put
\begin{equation}
 S=100M,\qquad
 \alpha=\inf_{0<z\le L}
 \frac{(2/\pi)\delta_*(M,z)-U_0(M,z)}{z}.
 \label{eq:parameter-alpha}
\end{equation}

\Needspace{5\baselineskip}
\begin{lemma}[Finiteness of the radius and positivity of the margin]
\label{lem:parameter-radius}
The set in \Cref{eq:parameter-radius} is nonempty.
Consequently \(M\) is a finite integer, and the number \(\alpha\)
in \Cref{eq:parameter-alpha} is strictly positive.
\end{lemma}

\begin{proof}
First, all sufficiently small \(t\ge0\) are regular.
The ground-level and symmetry conditions follow from
\Cref{prop:small-coupling}.  At zero coupling,
\[
 2\epsilon_b-E_3(0)=-\epsilon_e=\frac{k_0^2}{2}>0.49,
 \qquad 0.99<k_0<\gamma,
\]
where \Cref{prop:two-centre-spectrum} and the fixed value of \(D\)
give the strict numerical inequalities.  They persist for small
\(t\).  Positivity of the density and the direct-amplitude bound
are supplied by the uniform channel construction in
\Cref{prop:density-properties,lem:channel-expansion}.

We record why the fixed numbers in
\Cref{eq:fixed-parameter-numbers} are sufficient for the uniform
exterior tests.  Let \(\delta_0=\gamma-k_0>0\).
The estimate \Cref{eq:quantitative-density-drift} supplies one
remainder \(\rho_D(R)\downarrow0\), independent of all sufficiently
small \(t\), for the first three unnormalized expressions in
\Cref{eq:parameter-H}.  By \Cref{eq:two-centre-numerics},
\begin{equation}
 \delta_0<\frac{200}{D},\qquad
 100(\delta_0+\delta_0^2)<10^{-195}\ll\eta_0 .
 \label{eq:parameter-drift-slack}
\end{equation}
For the last line of \Cref{eq:parameter-H}, the uniform eventual
vertical ratio is at most \(20\), by \Cref{eq:strip-data}.
Consequently, for all sufficiently large fixed integers \(j\),
\[
 \begin{split}
 &\max\left\{H_0(j),\,\limsup_{t\downarrow0}H_t(j)\right\}\\
 &\qquad\le
 \max\left\{
   \frac{100(\delta_0+\delta_0^2)+\rho_D(j)}{\eta_0},
   \frac15\right\}<1.
 \end{split}
\]
The common remainder is obtained with \(D\) fixed, and the finite
radius \(j\) is then chosen before taking the small-coupling
limsup.  This uses uniform exterior estimates throughout.

We next verify the last condition in
\Cref{eq:parameter-radius}.  The scalar exterior expansions for
\(b,e\), including the simple odd factor of \(e\), imply that,
uniformly for \(0<z\le L\),
\[
 \partial_p\log(b/e)=-\delta_0+o(1),\qquad
 \partial_p^2\log(b/e)=o(1)\qquad(p\to\infty).
\]
The apparent singularity \(-\log z\) has no \(p\)-derivative.
For \(q>0\), both
\[
 \frac{q}{(1+q^2)\arctan q},
 \qquad
 \frac{q(1-q^2)}{(1+q^2)^2\arctan q}
\]
have absolute value at most \(1\); this follows from
\(q/(1+q^2)\le\arctan q\).
Applying these bounds to \(\delta_*=\arctan(\sqrt2 b/e)\)
therefore gives
\[
 \frac{|(\delta_*)_p|}{\delta_*}\le\delta_0+o(1),\qquad
 \frac{|(\delta_*)_{pp}|}{\delta_*}\le\delta_0^2+o(1).
\]
The bounds extend to \(z=0\) by continuity.  The cylindrical drift
coefficient satisfies \(\lvert p^{-1}+\partial_p\log n_0\rvert\le3\)
for sufficiently large \(p\).  Since
\(L_{n_0}\delta_*=q_*\delta_*\) with
\(0\le q_*\le2d_*\), it follows that
\[
 \frac{\partial_z(n_0\partial_z\delta_*)}{n_0\delta_*}
 \le 2d_*+3\delta_0+\delta_0^2+o(1)<\mu^2
\]
uniformly in the strip, once \(p\) is sufficiently large.
The last strict inequality uses
\(d_*<100/D\) and \Cref{eq:parameter-drift-slack}.
Thus the defining set for \(M\) is nonempty.

On the segment \(p=M\), let
\[
 f(z)=\frac2\pi\delta_*(M,z),\qquad
 v(z)=f(z)-U_0(M,z).
\]
The defining strict inequality for \(M\) gives
\[
 \frac1{n_0(M,z)}
 \partial_z\bigl(n_0(M,z)\partial_zv\bigr)-\mu^2v<0,
 \qquad
 v(0)=0,\quad v(L)>0.
\]
One-dimensional comparison and its boundary point conclusion
give \(v(z)>0\) for \(0<z\le L\) and \(v'(0)>0\).
The quotient \(v(z)/z\) consequently extends to a continuous
strictly positive function on the compact segment \([0,L]\).
Its minimum is positive, proving \(\alpha>0\).
\end{proof}

\subsection{The bad set and the charge}

For a regular \(t\), every folded minimizing angle for \(n_t\)
has a well-defined multiplier by \Cref{prop:angle-minimizer}.
In the following tests, the minimizing angle is not assumed to be
axisymmetric.  A supremum on the cylinder \(p=M\) includes all
azimuthal angles.

Define a subset \(\mathcal B\subset(0,1]\) as follows.
It contains every nonregular \(t\), every regular \(t\) with
\(H_t(M)\ge1\), and every remaining \(t\) for which \emph{there
exists a folded minimizing angle} \(\theta\), with multiplier \(d\),
that fails at least one of the tests
\begin{equation}
 \begin{gathered}
 0<d<0.001,\qquad
 \frac12\le\frac{n_t\cos^2\theta}{2b^2}\le2
       \quad\text{on } |x|\le2S,\\
 \sup_{\substack{p=M\\0\le z\le L}}
 \left(
   |\partial_z(\theta-\theta_*)|
   +|\partial_z(U_t-U_0)|
 \right)<\frac{\alpha}{2}.
 \end{gathered}
 \label{eq:parameter-angle-tests}
\end{equation}
Thus \(t\notin\mathcal B\) means that the tests hold for
\emph{every} folded minimizing angle.

\begin{definition}[The exact integer charge]
\label{def:explicit-charge}
Using only the preceding data for the fixed unit-charge problem,
define
\begin{equation}
 Z=2+\left\lceil
       \sup_{0<t\le1}\frac1t\,\mathbf1_{\mathcal B}(t)
      \right\rceil .
 \label{eq:explicit-charge}
\end{equation}
\end{definition}

\begin{lemma}[The charge is finite and passes all tests]
\label{lem:explicit-charge-finite}
The supremum in \Cref{eq:explicit-charge} is finite.  The resulting
\(Z\) is a positive integer, and \(t_*=1/Z\) is regular and lies
outside \(\mathcal B\).
\end{lemma}

\begin{proof}
By \Cref{lem:parameter-radius}, small couplings are regular and
the fixed-radius condition \(H_t(M)<1\) holds throughout some
interval \(0<t<\tau_1\).  We must check the tests for all minimizers,
not merely for a selected branch.

The convergence in \Cref{prop:angle-convergence} holds for every
sequence of folded minimizers as \(t\downarrow0\): their angles
converge locally uniformly to \(\theta_*\), their multipliers to
\(d_*\), and their derivatives converge uniformly on fixed exterior
compact sets.  The same statements are therefore uniform over
the collection of folded minimizers for sufficiently small \(t\);
otherwise a sequence violating uniform convergence would
contradict that proposition.  Since
\(0<d_*<100/D<0.001\), the multiplier test passes for all
minimizers.

The compact set \(\{|x|\le2S\}\) is now fixed.  Its reference
weight \(2b^2\) is bounded below by a positive constant.  Density
convergence and angle convergence give
\[
 \frac{n_t\cos^2\theta}{2b^2}\longrightarrow1
\]
uniformly there, for all folded minimizers.  This proves the
second test.  On the fixed cylinder \(p=M\), density convergence
also gives \(U_t(M,\cdot)\to U_0(M,\cdot)\) in \(C^1([0,L])\).
For example, this follows directly from the positive integral
series in \Cref{lem:strip-solution}: \(n_t\) converges uniformly
to a positive function on this compact segment, and the series
and its final-flux formula converge uniformly.  Together with
the \(C^1\) angle convergence and \(\alpha>0\), this proves the
last test.

Consequently \(\mathcal B\cap(0,\tau)=\varnothing\) for some
\(\tau>0\).  If
\[
 C_{\mathcal B}
 =\sup_{0<t\le1}\frac1t\,\mathbf1_{\mathcal B}(t),
\]
then \(0\le C_{\mathcal B}\le1/\tau\); in the empty-bad-set case
the supremum is zero.  Thus \(Z=2+\lceil C_{\mathcal B}\rceil\)
is finite and \(Z>C_{\mathcal B}\).  If \(t_*=1/Z\) belonged
to \(\mathcal B\), its contribution to the supremum would be
\(Z\), a contradiction.  Hence \(t_*\notin\mathcal B\), as claimed.
\end{proof}

\begin{proof}[Proof of \Cref{thm:main}]
Work first with the scaled system at \(t_*=1/Z>0\).
By \Cref{lem:explicit-charge-finite}, its regularity gives the
specified normalizable ground state in \(S_z=1/2\), and its
energy is an absolute ground energy.

It is useful to check explicitly that the fixed-parameter channel
theorem applies here.  We know that the selected coupling passes
the tests; we do not need to know that it belongs to the
particular perturbation interval used earlier.  Its ionic
isolation supplies a positive \(Q\)-gap.  The inequality
\(2\epsilon_b-E_3(t_*)>0.49\) permits a full-configuration
exponential weight with exponent \(a=0.98\), since
\(a^2/2=0.4802<0.49\).  On the cutoff region
\(|y|\ge r/3\), the full radius is at least
\(\sqrt{10/9}\,r\); hence, with room for the derivative estimates,
an exterior error exponent \(\sigma=1.02>\gamma\) is available.
Here \(\gamma<1.01\) follows from
\Cref{eq:two-centre-numerics}.  Moreover \(k_{t_*}<\gamma\)
makes the retained unperturbed fast term exponentially smaller
than \(A_{t_*}\) times any fixed inverse power.
Finally, \Cref{eq:regular-amplitude-test} supplies the nonzero
equatorial slope in the scalar channel coefficient.  The
complementary plane coefficient is nonzero because
\[
 Q(y_{1z}+y_{2z})\phi=(y_{1z}+y_{2z})\phi\ne0,
\]
by even ionic parity and the fact that a nonzero \(L^2\) vector
cannot be supported on the hyperplane \(y_{1z}+y_{2z}=0\).
The injective reduced inverse preserves nonzeroness.
These are exactly the fixed-parameter hypotheses of
\Cref{prop:fixed-parameter-channels}.  That proposition gives the
plane lower tail and both density translation limits at \(t_*\).

Choose any folded minimizing angle \(\theta\) for \(n_{t_*}\).
All tests in \Cref{eq:parameter-angle-tests} hold for it, and
\(H_{t_*}(M)<1\) supplies the drift and strip bounds with
\(\eta_0=10^{-60}\) and vertical ratio at most \(100\).
Also \(0.99<k_{t_*}<\gamma<1.01\) and \(M>10D\).
The hypotheses of \Cref{lem:quantitative-gap} are therefore
satisfied.  To indicate the numerical margins, its strip
derivative constant is below \(10^{30}\), so its horizontal
error is at most \(4\cdot10^{-30}U\), whereas
\(\mu^2-2d>0.008\).  At the inner cylinder, the reference gap
is at least \(\alpha z\); the derivative test reduces it by
less than \(\alpha z/2\).  The exterior gap is \(0.03\), while
the localization error \(10/S^2\) leaves a gap far larger
than \(d<0.001\).  Thus the full weighted gap holds for the
constructed orbitals \(g,h\).

The fixed-density obstruction in \Cref{thm:obstruction},
using the ensemble certificate and potential identification
from \Cref{sec:ensemble}, rules out every representing
potential in the stated class for this scaled density.
Finally, \Cref{prop:scaling} transfers the obstruction to the
physical system with charges \(Z,Z\).  The number \(Z\) is finite,
so the nuclear positions \((0,0,\pm D/Z)\) are finite and distinct.
This establishes \Cref{thm:main}.
\end{proof}

All quantities in the prescription are determined by the fixed
unit-charge auxiliary operators and their density-angle
minimizations.  Eigenvector phases do not enter, and the tests
quantify over every minimizing angle.  The supremum in
\Cref{eq:explicit-charge} is therefore an unambiguous definition
of an integer; it contains no test for representability by an
unknown potential.  Its nonnumerical character concerns the
evaluation of that integer, not the finiteness of the system or
the order of limits in the proof.
